\section{The Geneva Conventions and Additional Protocol I}
\label{sec:C.2}
This section is the translation of the four Geneva Conventions of 1949 and Additional Protocol I, in the form Table C.1 uses: for each article, the terms derived for a targeting deployment of Maven's kind, the terms considered and rejected and why, and the articles that yield none for this deployment. It exists because Chapter~\ref{sec:5} takes the set entire, and the law of armed conflict is the set that governs the motivating case. Most entries yield nothing: the Conventions mostly govern the treatment of persons in a party's hands, which a targeting pipeline doesn't do. The entries that yield terms are lifted into Table C.1. Articles the United States has rejected as customary are marked, and the terms derived from them are departures for the United States.

Every article of each instrument has an entry or belongs to one, in Table C.1's three parts: the article, the terms derived for a targeting deployment of Maven's kind, and the terms considered and rejected and why. An entry that yields no term says why in a few words.

\runin{Common Article 3}

\textbf{Common Art. 3 (conflicts not of an international character).} In a conflict not of an international character the person-directed terms apply as in an international one. Don't take part in violence to life or person, the taking of hostages, or sentences passed without a regularly constituted court, against persons taking no active part in hostilities, including those who have laid down their arms or are hors de combat. Rejected: apply only the rules of international conflict, because the floor reads both, and a state's description of its war as internal is its own claim.

The First Convention (wounded and sick in armed forces in the field)

\textbf{Arts. 1–2, 4–11 (general provisions).} None: scope, application and protecting powers; the floor applies the Conventions to every belligerent.

\textbf{Arts. 12–18 (the wounded and sick).} None: treatment of the wounded in a party's hands, which a targeting pipeline doesn't hold.

\textbf{Arts. 19–23 (medical units and establishments).} None separately: the medical-unit term derived from GC IV Art. 18 and AP I Art. 12 covers military medical units and establishments, and the deployment's translation says so. Rejected: a separate term for military units, because it would be redundant with that reading.

\textbf{Arts. 24–37 (medical personnel, buildings, transports).} None separately: the medical-unit term is read to include personnel, buildings and transports, including medical aircraft flying as agreed.

\textbf{Arts. 38–44 (the emblem).} None: misuse of the emblem is the executing unit's act, not the pipeline's.

\textbf{Arts. 45–64 (execution, repression of abuses, final provisions).} None: duties of states, not acts a pipeline performs.

The Second Convention (wounded, sick and shipwrecked at sea)

\textbf{Arts. 1–2, 4–11 (general provisions).} None: as in the First Convention.

\textbf{Arts. 12–21 (the wounded, sick and shipwrecked).} None: treatment at sea of persons in a party's hands.

\textbf{Arts. 22–40 (hospital ships, their personnel, medical transports).} None separately: the medical-unit term is read to include hospital ships, coastal rescue craft and medical transports at sea.

\textbf{Arts. 41–63 (the emblem, execution, repression, final provisions).} None: as in the First Convention.

The Third Convention (prisoners of war)

\textbf{Arts. 1–2, 4–16 (general provisions; general protection).} None: who is a prisoner of war and how prisoners are treated are custody questions.

\textbf{Arts. 17–108 (captivity: its beginning, internment, labour, finances, relations with the outside and with the authorities, discipline).} None: prisoner-of-war administration. A camp is a place where persons hors de combat are held, and the distinction terms already bar resolving one as a target.

\textbf{Arts. 109–143 (end of captivity, information bureaux, execution, final provisions).} None: repatriation and administration.

The Fourth Convention (civilians)

\textbf{Arts. 1–2, 4–12 (general provisions).} None: scope and application.

\textbf{Arts. 13–17 (field of application; hospital, safety and neutralized zones; evacuation).} None separately: zones established by agreement are protected places the medical-unit term and the distinction terms already cover.

\textbf{Arts. 18–19 (civilian hospitals; when protection ceases).} Don't resolve a location at a civilian hospital, or time an action to presence at one, absent a disinterested finding that it is used for acts harmful to the enemy, and never before the warning and time limit Art. 19 requires. (Lifted into Table C.1.) Rejected: lift on the operator's claim of misuse, because the claim is the operator's own.

\textbf{Arts. 20–22 (hospital staff; medical transport by land, sea and air).} None separately: within the medical-unit term.

\textbf{Art. 23 (free passage of relief consignments).} Don't take part in an action whose effect is to stop consignments the article requires a party to let through. (Within the siege term, Table C.1.) Rejected: don't take part in any interdiction of shipments, because the article lets a party check consignments and refuse them on stated grounds. The term reaches the effect on the population.

\textbf{Arts. 24–26 (children; family news; dispersed families).} None: welfare and family duties of a party, not acts a pipeline performs.

\textbf{Arts. 27–32 (treatment of protected persons; human shields; coercion; torture).} None separately: treatment in a party's hands. Art. 28's bar on using protected persons as shields binds the party that places them; for the attacker the action threshold weighs their presence. Rejected: treat a shielded objective as a lawful target without weighing the shields, because the floor counts them as people.

\textbf{Art. 33 (individual responsibility; collective penalties; pillage; reprisals).} Don't take part in an action whose stated purpose is to punish a population for the acts of others. (Lifted into Table C.1.)

\textbf{Art. 34 (hostages).} Don't take part in taking or holding hostages, or in locating people in order to take them.

\textbf{Arts. 35–46 (aliens in the territory of a party).} None: treatment of foreign nationals in a party's own territory. The removal terms (Table C.2) reach the analogous acts in peacetime.

\textbf{Arts. 47–48, 50–52 (occupied territory: inviolability of rights, repatriation, children, enlistment and labour).} None: duties of an occupying administration.

\textbf{Art. 49 (deportation and transfer).} Don't resolve locations, rank people or plan movements for a deportation or forcible transfer of protected persons out of occupied territory, or for the transfer of the occupier's own population into it. Rejected: don't take part in any evacuation, because Art. 49(2) permits an evacuation for the population's security or imperative military reasons, with return as soon as hostilities in the area end.

\textbf{Art. 53 (destruction of property in occupied territory).} None as a model term: the article's own exception, destruction rendered absolutely necessary by military operations, is a condition the deployment would find met on its own description (\S\ref{sec:18.3}, unconditional terms). Where the effect falls on survival, the siege term reaches it. Rejected: a term carrying the exception, because it would be found met by the party it constrains. Rejected: a term without it, because it would bar what the article permits.

\textbf{Art. 54 (judges and public officials).} None: occupation administration.

\textbf{Arts. 55 and 59 (food and medical supplies; relief schemes).} Within the siege term (Table C.1).

\textbf{Arts. 56–58, 60–78 (hygiene; hospitals; spiritual care; relief; penal legislation; security measures).} None: occupation administration and penal procedure.

\textbf{Arts. 79–141 (internment; information bureaux).} None: custody.

\textbf{Arts. 142–146, 148–159 (execution; final provisions).} None: duties of states.

\textbf{Art. 147 (grave breaches).} None new: each act it lists is refused by a term derived from its own article, and the list adds no act a targeting pipeline performs beyond them. Rejected: a separate grave-breaches term, because it would be redundant.

\runin{Additional Protocol I}

\textbf{Art. 1(1)–(3), Arts. 2–7 (general provisions).} None: scope, definitions and supervision.

\textbf{Art. 1(4) (wars against colonial domination, alien occupation and racist regimes).} None of its own: it treats such conflicts as international, which is what the occupation exception's finding does for the force term (\S\ref{sec:6.7}; Appendix~\ref{sec:F.5}). Marked: the United States rejects it as customary.

\textbf{Arts. 8–11 (definitions; scope; protection and care; protection of persons).} None: the wounded and sick, and protection of persons in a party's hands.

\textbf{Art. 12 (protection of medical units).} With GC IV Art. 18, the medical-unit term (Table C.1).

\textbf{Arts. 13–34 (end of protection; medical personnel; medical transports; the missing and the dead).} None separately: the medical-unit term is read to include civilian medical personnel and medical transports by land, sea and air; the missing and the dead are a party's duties.

\textbf{Art. 35 (basic rules on methods and means).} None of its own: its first paragraph is the principle the specific terms carry out, and its second and third concern the effects of weapons the pipeline doesn't choose.

\textbf{Art. 36 (new weapons).} None: a duty of review before adoption. \S\ref{sec:7.1} notes it as where the review ratio's parameters would naturally be assessed.

\textbf{Art. 37 (perfidy).} None for this deployment: perfidy is committed by the unit that feigns protected status, at the handoff, not by a pipeline that resolves locations. A deployment that plans the feigning yields the term.

\textbf{Arts. 38–39 (recognized emblems; emblems of nationality).} None: the executing unit's acts.

\textbf{Arts. 40–42 (quarter; safeguard of an enemy hors de combat; occupants of aircraft).} None separately: an order of no quarter is given to the executing force, and a person who surrenders after release is the handoff's and the cancellation window's case (Art. 57(2)(b)). Recognizing surrender is the first case a system without a reviewer has to pass (\S\ref{sec:7.1}).

\textbf{Arts. 43, 45–47 (armed forces; persons who took part in hostilities; spies; mercenaries).} None: status determinations made for treatment in custody.

\textbf{Art. 44 (combatants and prisoners of war).} None: it relaxes the duty of distinction for combatant status, a question about who is owed prisoner-of-war treatment, which no term in this floor administers. Marked: the United States rejects it as customary.

\textbf{Art. 48 (basic rule of distinction).} With Art. 52 and the Covenant's Art. 3, the person-directed terms (the machine-selection departure; Table C.1).

\textbf{Arts. 49–50 (definitions of attacks and of civilians).} None separately: the terms use the definitions. Art. 50(1)'s presumption that a person whose status is in doubt is a civilian is what the action threshold carries.

\textbf{Art. 51 (protection of the civilian population).} Paragraphs (4)–(5) yield the action threshold and the model term in Table C.1; paragraph (6) yields the collective-punishment term with GC IV Art. 33. Paragraph (2) joins paragraph (6) in that entry: don't take part in an action whose stated purpose is to spread terror among the civilian population. (Lifted into Table C.1.) Paragraph (7) binds the party that places civilians as shields.

\textbf{Art. 52 (protection of civilian objects).} Within the person-directed terms (the machine-selection departure; Table C.1).

\textbf{Art. 53 (cultural objects and places of worship).} Don't resolve a location at a historic monument, work of art or place of worship that constitutes the cultural or spiritual heritage of peoples, or time an action to presence at one, absent a disinterested finding that it is used in support of the military effort. (Lifted into Table C.1.) Rejected: lift on the operator's claim of military use, because the claim is the operator's own.

\textbf{Art. 54 (objects indispensable to survival).} The siege term (Table C.1).

\textbf{Art. 55 (the natural environment).} None for this deployment: it governs means expected to cause widespread, long-term and severe damage, which a pipeline that resolves locations doesn't choose; the dangerous-forces term carries the nearest case.

\textbf{Art. 56 (works containing dangerous forces).} The dangerous-forces term (Table C.1).

\textbf{Art. 57 (precautions in attack).} The review ratio's home: verification (57(2)(a)(i)), proportionality (57(2)(a)(iii)) and cancellation (57(2)(b)) are what the operator term pre-registers. The deployment states how it gives the effective advance warning of 57(2)(c) where circumstances permit. Rejected: a model-side term requiring review, because the model can't establish whether review occurred.

\textbf{Art. 58 (precautions against the effects of attacks).} None for this deployment: it binds the defending party, which must remove civilians from the vicinity of military objectives and avoid siting objectives in populated areas. A defender's pipeline that plans positions yields the term.

\textbf{Arts. 59–60 (non-defended localities; demilitarized zones).} Don't resolve a location in a locality declared non-defended or a zone agreed demilitarized, unless a disinterested body finds the article's conditions no longer met. (Lifted into Table C.1.) Rejected: lift on the operator's claim that the locality no longer meets the conditions, because the claim is the operator's own.

\textbf{Arts. 61–67 (civil defence).} Don't resolve a location at a civil-defence organization, its personnel, buildings or shelters, absent a disinterested finding that they are used for acts harmful to the enemy, and after the warning Art. 65 requires. (Within the medical-unit entry of Table C.1.)

\textbf{Arts. 68–71 (relief in favour of the civilian population).} Within the siege term.

\textbf{Arts. 72–79 (treatment of persons in a party's power; refugees; families; fundamental guarantees; women; children; journalists).} None separately: treatment in custody; journalists are civilians, and the distinction terms cover them.

\textbf{Arts. 80–84, 86–91 (execution; failure to act; duty of commanders; mutual assistance; the Fact-Finding Commission).} None: duties of states and commanders. The record the review ratio lodges is evidence of what a commander knew under Arts. 86(2) and 87.

\textbf{Art. 85 (grave breaches).} None new: as with GC IV Art. 147.

\textbf{Arts. 92–102 (final provisions).} None.
